在过去几年里，将 Transformer 扩展到更长的序列长度一直是一个关键问题，这有望提升语言建模与高分辨率图像理解性能，并解锁代码、音频与视频生成等新应用。
注意力层是扩展到长序列时的主要瓶颈，因为其运行时间和内存开销都随序列长度呈二次增长。
\sysnameone~\citep{dao2022flashattention} 利用了 GPU 非对称的内存层次结构，在不做近似的前提下，实现了显著的内存节省（由二次降为线性）和运行加速（相对优化基线提升 2-4$\times$）。
然而，\sysnameone 相比优化后的矩阵乘（GEMM）仍然不够快，只能达到理论最大 FLOPs/s 的 25-40\%。
我们观察到，这种低效率来自 GPU 上不同 thread block 与 warp 之间次优的工作划分，导致要么占用率低，要么发生不必要的共享内存读写。
我们提出 \sysname，通过更优的工作划分来解决这些问题。
具体而言，我们 (1) 调整算法以减少非矩阵乘 FLOPs 数量；(2) 即使对单个头，也在不同 thread block 之间并行注意力计算以提高占用率；(3) 在每个 thread block 内，在不同 warp 之间分配工作以减少通过共享内存的通信。
这些改进相较 \sysnameone 带来约 2$\times$ 加速，在 A100 上达到理论最大 FLOPs/s 的 50-73\%，并接近 GEMM 操作的效率。
我们进一步在端到端 GPT 风格模型训练中实证验证：\sysname 在每张 A100 GPU 上可达到最高 225 TFLOPs/s 的训练速度（模型 FLOPs 利用率 72\%）。\footnote{\sysname
  代码见 \url{https://github.com/Dao-AILab/flash-attention}}

% 与 \sysnameone 相比，在相同时间下可支持最高 2$\times$ 的序列长度，
% 并带来更好的下游性能。\footnote{\sysname
%   代码见 \url{https://github.com/Dao-AILab/flash-attention}}
